\documentclass{amsart}
\usepackage[UTF8,fontset=windows]{ctex}
\usepackage[a4paper,margin=24mm]{geometry}
% For dealing with references we use the comment environment
\usepackage{verbatim}
\newenvironment{reference}{\comment}{\endcomment}
%\newenvironment{reference}{}{}
\newenvironment{slogan}{\comment}{\endcomment}
\newenvironment{history}{\comment}{\endcomment}

% For commutative diagrams we use Xy-pic
\usepackage[all]{xy}

% We use 2cell for 2-commutative diagrams.
\xyoption{2cell}
\UseAllTwocells

% We use multicol for the list of chapters between chapters
\usepackage{multicol}

% This is generally recommended for better output
\usepackage{lmodern}
\usepackage[T1]{fontenc}

% For cross-file-references

% Package for hypertext links:
\usepackage{hyperref}

% For any local file, say "hello.tex" you want to link to please

% Theorem environments.
%
\theoremstyle{plain}
\newtheorem{theorem}[subsection]{Theorem}
\newtheorem{proposition}[subsection]{Proposition}
\newtheorem{lemma}[subsection]{Lemma}

\theoremstyle{definition}
\newtheorem{definition}[subsection]{Definition}
\newtheorem{example}[subsection]{Example}
\newtheorem{exercise}[subsection]{Exercise}
\newtheorem{situation}[subsection]{Situation}

\theoremstyle{remark}
\newtheorem{remark}[subsection]{Remark}
\newtheorem{remarks}[subsection]{Remarks}

\numberwithin{equation}{subsection}

% Macros
%
\def\lim{\mathop{\mathrm{lim}}\nolimits}
\def\colim{\mathop{\mathrm{colim}}\nolimits}
\def\Spec{\mathop{\mathrm{Spec}}}
\def\Hom{\mathop{\mathrm{Hom}}\nolimits}
\def\Ext{\mathop{\mathrm{Ext}}\nolimits}
\def\SheafHom{\mathop{\mathcal{H}\!\mathit{om}}\nolimits}
\def\SheafExt{\mathop{\mathcal{E}\!\mathit{xt}}\nolimits}
\def\Sch{\mathit{Sch}}
\def\Mor{\mathop{\mathrm{Mor}}\nolimits}
\def\Ob{\mathop{\mathrm{Ob}}\nolimits}
\def\Sh{\mathop{\mathit{Sh}}\nolimits}
\def\NL{\mathop{N\!L}\nolimits}
\def\CH{\mathop{\mathrm{CH}}\nolimits}
\def\proetale{{pro\text{-}\acute{e}tale}}
\def\etale{{\acute{e}tale}}
\def\QCoh{\mathit{QCoh}}
\def\Ker{\mathop{\mathrm{Ker}}}
\def\Im{\mathop{\mathrm{Im}}}
\def\Coker{\mathop{\mathrm{Coker}}}
\def\Coim{\mathop{\mathrm{Coim}}}

% Boxtimes
%
\DeclareMathSymbol{\boxtimes}{\mathbin}{AMSa}{"02}

%
% Macros for moduli stacks/spaces
%
\def\QCohstack{\mathcal{QC}\!\mathit{oh}}
\def\Cohstack{\mathcal{C}\!\mathit{oh}}
\def\Spacesstack{\mathcal{S}\!\mathit{paces}}
\def\Quotfunctor{\mathrm{Quot}}
\def\Hilbfunctor{\mathrm{Hilb}}
\def\Curvesstack{\mathcal{C}\!\mathit{urves}}
\def\Polarizedstack{\mathcal{P}\!\mathit{olarized}}
\def\Complexesstack{\mathcal{C}\!\mathit{omplexes}}
% \Pic is the operator that assigns to X its picard group, usage \Pic(X)
% \Picardstack_{X/B} denotes the Picard stack of X over B
% \Picardfunctor_{X/B} denotes the Picard functor of X over B
\def\Pic{\mathop{\mathrm{Pic}}\nolimits}
\def\Picardstack{\mathcal{P}\!\mathit{ic}}
\def\Picardfunctor{\mathrm{Pic}}
\def\Deformationcategory{\mathcal{D}\!\mathit{ef}}

\hypersetup{hidelinks}
\widowpenalty=10000\clubpenalty=10000\displaywidowpenalty=10000
\begin{document}
\title[Stacks Project, Tag 0CDA]{亏格零曲线极小模型非唯一时的特殊纤维}
\maketitle
\noindent\textbf{Counting scope.} The primary problem is Tag 0CDA, marked below. All preceding numbered statements are supporting author context and are not counted as additional problems. This is a local pilot with a temporary ID, not production admission.


\noindent Copyright (C) 2005--2025 Johan de Jong. Stacks Project authors. Source: \href{https://stacks.math.columbia.edu/tag/0CDA}{Tag 0CDA}.

\noindent Difficulty H2: 本题从亏格零曲线存在两个不同极小模型推出每个特殊纤维都是射影直线。证明比较两个模型的共同爆破，先收缩特殊纤维重数大于一的异常曲线降低爆破次数，再分析未收缩的异常曲线与原特殊纤维的同构。允许曲面双有理分解和异常曲线收缩后，重数分析、归纳及模型比较仍由本证明完成，判H2。.

\noindent History: excerpt from revision local cached snapshot; exact chapter hashes; revision unknown, models.tex, lines 5543--5649. Only reference presentation was adapted. Licensed under GNU FDL 1.2 or later; no invariant sections or cover texts. See accompanying COPYING.

\section{Models}
\label{section-models}

\noindent
In this chapter $R$ will be a discrete valuation ring and $K$ will
be its fraction field. If needed we will denote $\pi \in R$ a
uniformizer and $k = R/(\pi)$ its residue field.

\medskip\noindent
Let $V$ be an algebraic $K$-scheme
(Varieties, Definition \ref{varieties-definition-algebraic-scheme}).
A {\it model} for $V$ will
mean a flat finite type\footnote{Occasionally it is useful to
allow models to be locally of finite type over $R$, but we'll
cross that bridge when we come to it.}
morphism $X \to \Spec(R)$ endowed with
an isomorphism $V \to X_K = X \times_{\Spec(R)} \Spec(K)$. We often
will identify $V$ and the generic fibre $X_K$ of $X$ and
just write $V = X_K$.
The special fibre is $X_k = X \times_{\Spec(R)} \Spec(k)$.
A {\it morphism of models $X \to X'$ for $V$}
is a morphism $X \to X'$ of schemes over $R$ which induces
the identity on $V$.

\medskip\noindent
We will say {\it $X$ is a proper model of $V$} if $X$ is a model
of $V$ and the structure morphism $X \to \Spec(R)$ is proper.
Similarly for separated models, smooth models, and add more here.
We will say {\it $X$ is a regular model of $V$} if $X$ is a model
of $V$ and $X$ is a regular scheme.
Similarly for normal models, reduced models, and add more here.

\medskip\noindent
Let $R \subset R'$ be an extension of discrete valuation rings
(More on Algebra, Definition
\ref{more-algebra-definition-extension-discrete-valuation-rings}).
This induces an extension $K'/K$ of fraction fields.
Given an algebraic scheme $V$ over $K$, denote $V'$ the
base change $V \times_{\Spec(K)} \Spec(K')$. Then there is
a functor
$$
\text{models for }V\text{ over }R
\longrightarrow
\text{models for }V'\text{ over }R'
$$
sending $X$ to $X \times_{\Spec(R)} \Spec(R')$.



\begin{definition}
\label{definition-minimal-model}
Let $C$ be a smooth projective curve over $K$ with
$H^0(C, \mathcal{O}_C) = K$. A {\it minimal model}
will be a regular, proper model $X$ for $C$ such that
$X$ does not contain an exceptional curve of the first kind
(Resolution of Surfaces, Section \href{https://stacks.math.columbia.edu/tag/0C2I}{section minus one (Tag 0C2I)}).
\end{definition}

\begin{definition}
\label{definition-numerical-type-model}
In Situation \ref{situation-regular-model} the
{\it numerical type associated to $X$} is the numerical
type described in Lemma \href{https://stacks.math.columbia.edu/tag/0CA4}{lemma numerical type of model (Tag 0CA4)}.
\end{definition}

\begin{lemma}
\label{lemma-numerical-type-minimal-model}
In Situation \ref{situation-regular-model}. The following
are equivalent
\begin{enumerate}
\item $X$ is a minimal model, and
\item the numerical type associated to $X$ is minimal.
\end{enumerate}
\end{lemma}

\begin{lemma}[Genus zero]
\label{lemma-genus-zero}
The only minimal numerical type of genus zero is
$n = 1$, $m_1 = 1$, $a_{11} = 0$, $w_1 = 1$, $g_1 = 0$.
\end{lemma}

\begin{proof}
Follows from Lemmas \href{https://stacks.math.columbia.edu/tag/0C7B}{lemma non irreducible minimal type genus at least one (Tag 0C7B)}
and \href{https://stacks.math.columbia.edu/tag/0C73}{lemma irreducible (Tag 0C73)}.
\end{proof}


\begin{definition}
\label{varieties-definition-algebraic-scheme}
Let $k$ be a field. An {\it algebraic $k$-scheme} is a scheme $X$ over $k$
such that the structure morphism $X \to \Spec(k)$ is of
finite type. A {\it locally algebraic $k$-scheme} is a scheme $X$ over $k$
such that the structure morphism $X \to \Spec(k)$ is
locally of finite type.
\end{definition}

\begin{definition}
\label{more-algebra-definition-extension-discrete-valuation-rings}
We say that $A \to B$ or $A \subset B$ is an
{\it extension of discrete valuation rings} if $A$ and $B$ are
discrete valuation rings and $A \to B$ is injective and local.
In particular, if $\pi_A$ and $\pi_B$ are uniformizers of
$A$ and $B$, then $\pi_A = u \pi_B^e$ for some $e \geq 1$ and unit
$u$ of $B$. The integer $e$ does not depend on the choice of
the uniformizers as it is also the unique integer $\geq 1$ such that
$$
\mathfrak m_A B = \mathfrak m_B^e
$$
The integer $e$ is called the {\it ramification index} of $B$ over $A$.
We say that $B$ is {\it weakly unramified} over $A$ if $e = 1$.
If the extension of residue fields
$\kappa_A = A/\mathfrak m_A \subset \kappa_B = B/\mathfrak m_B$
is finite, then we set $f = [\kappa_B : \kappa_A]$ and we
call it the {\it residual degree} or {\it residue degree}
of the extension $A \subset B$.
\end{definition}

\begin{situation}
\label{situation-regular-model}
Let $R$ be a discrete valuation ring with fraction field $K$,
residue field $k$, and uniformizer $\pi$.
Let $C$ be a smooth projective curve over $K$ with $H^0(C, \mathcal{O}_C) = K$.
Let $X$ be a regular proper model of $C$.
Let $C_1, \ldots, C_n$ be the irreducible components of the special
fibre $X_k$. Write $X_k = \sum m_i C_i$ as in
Lemma \href{https://stacks.math.columbia.edu/tag/0C5Z}{lemma divisor special fiber (Tag 0C5Z)}.
\end{situation}

\section*{Primary problem: Tag 0CDA}
\begin{lemma}
\label{lemma-nonuniqueness}
Let $C$ be a smooth projective curve over $K$ with $H^0(C, \mathcal{O}_C) = K$
and genus $0$. If there is more than one minimal model for $C$, then
the special fibre of every minimal model is isomorphic to $\mathbf{P}^1_k$.
\end{lemma}

\noindent
This lemma can be improved to say that the birational transformation
between two nonisomorphic minimal models can be factored as a sequence of
elementary transformations as in Example \href{https://stacks.math.columbia.edu/tag/0CA0}{example nonunique in genus zero (Tag 0CA0)}.
If we ever need this, we will precisely formulate and prove this here.

\begin{proof}
Let $X$ be some minimal model of $C$. The numerical type associated to $X$
has genus $0$ and is minimal (Definition \ref{definition-numerical-type-model}
and Lemma \ref{lemma-numerical-type-minimal-model}).
Hence by Lemma \ref{lemma-genus-zero} we see that
$X_k$ is reduced, irreducible, has $H^0(X_k, \mathcal{O}_{X_k}) = k$,
and has genus $0$. Let $Y$ be a second minimal model for $C$
which is not isomorphic to $X$. By
Resolution of Surfaces, Lemma \href{https://stacks.math.columbia.edu/tag/0C5S}{lemma birational regular surfaces (Tag 0C5S)}
there exists a diagram of $S$-morphisms
$$
X = X_0 \leftarrow X_1 \leftarrow \ldots \leftarrow X_n = Y_m
\to \ldots \to Y_1 \to Y_0 = Y
$$
where each morphism is a blowup in a closed point. We will prove the
lemma by induction on $m$. The base case is $m = 0$; it is
true in this case because we assumed that $Y$ is minimal
hence this would mean $n = 0$, but $X$ is not isomorphic
to $Y$, so this does not happen, i.e., there is nothing to check.

\medskip\noindent
Before we continue, note that $n + 1 = m + 1$ is equal to the number of
irreducible components of the special fibre of $X_n = Y_m$ because
both $X_k$ and $Y_k$ are irreducible. Another observation we
will use below is that if $X' \to X''$ is a morphism
of regular proper models for $C$, then $X' \to X''$ is an isomorphism
over an open set of $X''$ whose complement is a finite set
of closed points of the special fibre of $X''$, see
Varieties, Lemma
\href{https://stacks.math.columbia.edu/tag/0BFP}{lemma modification normal iso over codimension 1 (Tag 0BFP)}.
In fact, any such $X' \to X''$ is a sequence
of blowing ups in closed points (Resolution of Surfaces, Lemma
\href{https://stacks.math.columbia.edu/tag/0C5R}{lemma proper birational regular surfaces (Tag 0C5R)}) and the
number of blowups is the difference in the number of irreducible
components of the special fibres of $X'$ and $X''$.

\medskip\noindent
Let $E_i \subset Y_i$, $m \geq i \geq 1$ be the curve which is contracted
by the morphism $Y_i \to Y_{i - 1}$. Let $i$ be the biggest index such
that $E_i$ has multiplicity $> 1$ in the special fibre of $Y_i$.
Then the further blowups $Y_m \to \ldots \to Y_{i + 1} \to Y_i$
are isomorphisms over $E_i$ since otherwise $E_j$ for some $j > i$
would have multiplicity $> 1$. Let $E \subset Y_m$ be the inverse
image of $E_i$. By what we just said $E \subset Y_m$ is an
exceptional curve of the first kind. Let $Y_m \to Y'$ be the
contraction of $E$ (which exists by Resolution of Surfaces, Lemma
\href{https://stacks.math.columbia.edu/tag/0C2N}{lemma contract when quasi projective (Tag 0C2N)}). The morphism
$Y_m \to X$ has to contract $E$, because $X_k$ is reduced.
Hence there are morphisms $Y' \to Y$ and $Y' \to X$ (by
Resolution of Surfaces, Lemma
\href{https://stacks.math.columbia.edu/tag/0C5J}{lemma factor through contraction (Tag 0C5J)})
which are compositions of at most $n - 1 = m - 1$ contractions
of exceptional curves (see discussion above). We win by induction on $m$.
Upshot: we may assume that the special fibres of all of
the curves $X_i$ and $Y_i$ are reduced.

\medskip\noindent
Since the fibres of $X_i$ and $Y_i$ are reduced, it has to be the
case that the blowups $X_i \to X_{i - 1}$ and $Y_i \to Y_{i - 1}$
happen in closed points which are regular points of the special fibres.
Namely, if $X''$ is a regular model for $C$ and if $x \in X''$
is a closed point of the special fibre, and
$\pi \in \mathfrak m_x^2$, then the exceptional fibre $E$ of the
blowup $X' \to X''$ at $x$ has multiplicity at least $2$ in the
special fibre of $X'$ (local computation omitted).
Hence $\mathcal{O}_{X''_k, x} = \mathcal{O}_{X'', x}/\pi$ is
regular (Algebra, Lemma \href{https://stacks.math.columbia.edu/tag/00NQ}{lemma regular ring CM (Tag 00NQ)}) as claimed.
In particular $x$ is a Cartier divisor on the unique
irreducible component $Z'$ of $X''_k$ it lies on
(Varieties, Lemma \href{https://stacks.math.columbia.edu/tag/0B8Y}{lemma regular point on curve (Tag 0B8Y)}).
It follows that the strict transform $Z \subset X'$ of $Z'$
maps isomorphically to $Z'$
(use Divisors, Lemmas \href{https://stacks.math.columbia.edu/tag/080E}{lemma strict transform (Tag 080E)} and
\href{https://stacks.math.columbia.edu/tag/0807}{lemma blow up effective Cartier divisor (Tag 0807)}).
In other words, if an irreducible component $Z$ of $X_i$
is not contracted under the map $X_i \to X_j$ ($i > j$)
then it maps isomorphically to its image.

\medskip\noindent
Now we are ready to prove the lemma.
Let $E \subset Y_m$ be the exceptional curve of the first kind
which is contracted by the morphism $Y_m \to Y_{m - 1}$. If $E$ is
contracted by the morphism $Y_m = X_n \to X$, then there is a factorization
$Y_{m - 1} \to X$ (Resolution of Surfaces, Lemma
\href{https://stacks.math.columbia.edu/tag/0C5J}{lemma factor through contraction (Tag 0C5J)})
and moreover $Y_{m - 1} \to X$ is a sequence of blowups
in closed points (Resolution of Surfaces, Lemma
\href{https://stacks.math.columbia.edu/tag/0C5R}{lemma proper birational regular surfaces (Tag 0C5R)}).
In this case we lower $m$ and we win by induction.
Finally, assume that $E$ is not contracted by the morphism $Y_m \to X$.
Then $E \to X_k$ is surjective as $X_k$ is irreducible
and by the above this means it is an isomorphism.
Hence $X_k$ is isomorphic to a projective line as desired.
\end{proof}
\end{document}
